% GENERATED, edit the French unit or the translation table
\mxnewpage\section*{EV4-T1 -- EAM preparation}\label{EV4-T1}
\addcontentsline{toc}{section}{EV4-T1 -- EAM preparation}
\textbf{Practice paper}. Whole programme.\par Duration: 2 h. No calculator. Marked out of 20. Practice sessions in May and June; EAM date for the zone to be confirmed.
\subsection*{Part 1 -- Fluency (6 points)}
Twelve questions, 0.5 point each. Only one correct answer; no justification required. An incorrect or missing answer scores 0, with no loss of points.
\begin{enumerate}
\item Calculate $5+2\times3$.\par A. $14$\quad B. $11$\quad C. $21$\quad D. $10$
\item Calculate $\frac12+\frac13$.\par A. $\frac{5}{3}$\quad B. $\frac{5}{6}$\quad C. $\frac{1}{6}$\quad D. $\frac{2}{5}$
\item Calculate $\dfrac{2^5}{2^{3}}$.\par A. $1$\quad B. $4$\quad C. $10$\quad D. $2$
\item A value of 50 increases by 20 \%. Final value?\par A. $40$\quad B. $10$\quad C. $60$\quad D. $70$
\item A value of 80 is obtained after a 20 \% decrease. Initial value?\par A. $60$\quad B. $96$\quad C. $100$\quad D. $64$
\item A 25 \% increase then a 20 \% decrease: overall coefficient?\par A. $\frac{4}{5}$\quad B. $1$\quad C. $\frac{9}{10}$\quad D. $\frac{21}{20}$
\item The solution of $2x+5=0$ is\par A. $\frac{5}{2}$\quad B. $10$\quad C. $\frac{-5}{2}$\quad D. $-10$
\item The slope between $(0;5)$ and $(2;11)$ is\par A. $3$\quad B. $-3$\quad C. $6$\quad D. $2$
\item The mean of $5, 7, 9$ is\par A. $7$\quad B. $21$\quad C. $5$\quad D. $8$
\item A sector of $90^\circ$ represents what fraction of the total?\par A. $\frac{1}{4}$\quad B. $\frac{3}{4}$\quad C. $\frac{1}{3}$\quad D. $\frac{1}{2}$
\item A probability is $\dfrac{2}{5}$. Its complement is\par A. $\frac{3}{5}$\quad B. $\frac{1}{5}$\quad C. $\frac{2}{5}$\quad D. $\frac{5}{2}$
\item For $f(x)=x^2+5$, $f(-2)$ is\par A. $20$\quad B. $7$\quad C. $1$\quad D. $9$
\end{enumerate}
\mxnewpage\subsection*{Part 2 -- Independent exercises (14 points)}
Every step must be justified. Exact results are expected unless stated otherwise.
\mxneed{12}
\subsection*{Exercise 1 -- Choosing a distribution (5 points)}
We compare two games. The first gives a gain $X=2$ with certainty. The second gives a gain $Y=-1$ with probability $\dfrac23$ and a gain $Y=8$ with probability $\dfrac13$. The gains are expressed in the same monetary unit; a negative value represents a loss.
\begin{enumerate}
\item Calculate $E(X)$ and $E(Y)$. \hfill(1 pt)
\item Calculate $V(Y)$ and compare the dispersions of the two gains. \hfill(2 pts)
\item Both games are played and the gains added. Calculate $E(X+Y)$, citing the property used. \hfill(2 pts)
\end{enumerate}
\mxneed{12}
\subsection*{Exercise 2 -- Discrete threshold and continuous model (5 points)}
A quantity is modelled by $Q(t)=64e^{-t}$ for $t\geq0$. It is observed at integer times $n\in\N$ and we set $u_n=Q(n)$. For comparisons without a calculator, the bounds $e^2<8<e^3$ are given.
\begin{enumerate}
\item Study the direction of variation. \hfill(1 pt)
\item Determine the first integer n such that $Q(n)<8$ using the bounds provided. \hfill(2 pts)
\item Give the type of $u_n=Q(n)$ and the expression of $u_0+\cdots+u_3$. \hfill(2 pts)
\end{enumerate}
\mxneed{12}
\subsection*{Exercise 3 -- Equation and circle (4 points)}
In an orthonormal frame, consider the set $\mathcal C$ defined by $x^2+y^2-2x-4y+1=0$. The task is to determine its characteristics, then to identify a particular tangent.
\begin{enumerate}
\item Show that $\mathcal C$ is a circle and state its centre and radius. \hfill(1 pt)
\item Determine its points of intersection with the line with equation $x=1$. \hfill(1 pt)
\item Determine the equation of the tangent to the circle at the point with the greatest y-coordinate and justify. \hfill(2 pts)
\end{enumerate}
